\selectlanguage{american}%

\section{Ordinary Differential Equations}

We consider the first-order ODE 
\begin{align*}
\frac{\d}{\d t}x(t)= & ~F(x(t),t),\quad\forall0\leq t\leq T,\\
x(0)= & ~v.
\end{align*}
where $F:\R^{d+1}\rightarrow\R$. The basic idea to solve it comes
from the Picard-Lindel�f theorem, a constructive existence proof for
a large class of ODEs.

\subsection{Picard-Lindel�f theorem}

\label{sec:picard_lindelof} In general, we can rewrite the first-order
ODE as an integral equation 
\[
x(t)=v+\int_{0}^{t}F(x(s),s)\d s\quad\text{for all }0\leq t\leq T.
\]

To simplify the notation, we use $\mathcal{C}([0,T],\R^{d})$ to denote
$\R^{d}$-valued functions on $[0,T]$. We define the operator $\T$
from $\mathcal{C}([0,T],\R^{d})$ to $\mathcal{C}([0,T],\R^{d})$
by 
\begin{equation}
\T(x)(t)=v+\int_{0}^{t}F(x(s),s)\d s\quad\text{for all }0\leq t\leq T.\label{eq:def_T_operator}
\end{equation}
Therefore, the integral equation is simply $x=\T(x)$. Let $\T^{\circ j}$
be the composition of $j$ many $\T$, i.e., 
\begin{align*}
\T^{\circ j}(x)=\underbrace{\T(\T(\cdots\T}_{j~\text{many}~\T}(x)\cdots)).
\end{align*}

The Banach fixed point theorem says that the integral equation $x=\T(x)$
has a unique solution if there is a norm, and $j\in\mathbb{N}$ such
that the map $\T^{\circ j}$ has Lipschitz constant less than 1.

The Picard-Lindel�f theorem shows that if $F$ is Lipschitz in $x$,
then the map $\T^{\circ j}$ has Lipschitz constant less than 1 for
some positive integer $j$.
\begin{thm}
\label{lem:Lipschitz_T} Given any norm $\|\cdot\|$ on $\R^{d}$.
Let $L$ be the Lipschitz constant of $F$ in $x$, i.e. 
\[
\|F(x,s)-F(y,s)\|\leq L\|x-y\|\quad\text{for all }x,y\in\R^{d},s\in[0,T].
\]
For any $x\in\mathcal{C}([0,T],\R^{d})$, we define the corresponding
norm 
\[
\|x\|\defeq\max_{0\leq t\leq T}\|x(t)\|.
\]
Then, the Lipschitz constant of $\T^{\circ j}$ in this norm is upper
bounded by $(LT)^{j}/j!$.
\end{thm}

\begin{proof}
We prove this lemma with a stronger induction hypothesis
\[
\|(\T^{\circ j}x)(h)-(\T^{\circ j}y)(h)\|\leq\frac{L^{j}h^{j}}{j!}\|x-y\|\quad\text{for all }0\leq h\leq T.
\]
The base case $j=0$ is trivial. For the inductive step, 
\begin{align*}
~\|\T^{\circ j}x(h)-\T^{\circ j}y(h)\|= & ~\|\T\T^{\circ(j-1)}x(h)-\T\T^{\circ(j-1)}y(h)\|\\
%=\leq
\end{align*}
This completes the proof.
\end{proof}
To make the Picard-Lindel�f theorem algorithmic, we have to decide
how to represent a function in $\mathcal{C}([0,T],\R^{d})$. One standard
way is to use a polynomial $p_{i}(t)$ in $t$ for each coordinate
$i\in[d]$. More generally, suppose there is a basis $\{\varphi_{j}\}_{j=1}^{D}\subset\mathcal{C}([0,T],\R)$
such that for all $i\in[d]$, $\frac{\d x_{i}}{\d t}$ is approximated
by some linear combination of $\varphi_{j}(t)$. For example, if $\varphi_{j}(t)=t^{j-1}$
for $j\in[D]$, then our assumption is simply saying $\frac{\d x_{i}}{\d t}$
is approximated by a degree $D-1$ polynomial. Other possible basis
are piecewise-polynomial and Fourier series. By Gram-Schmidt orthogonalization,
we can always pick points $\{c_{i}\}_{j=1}^{D}$ such that 
\[
\varphi_{j}(c_{i})=\delta_{i,j}\quad\text{for }i,j\in[D].
\]
The benefit of such a basis is that for any $f\in\text{span}(\varphi_{j})$,
we have that $f(t)=\sum_{j=1}^{D}f(c_{j})\varphi_{j}(t).$

For polynomials, we can use the basis consisting of the Lagrange polynomials:
\[
\varphi_{j}(t)=\prod_{i\in[D]\backslash\{j\}}\frac{t-c_{i}}{c_{j}-c_{i}}\quad\text{for }j\in[D].
\]
The only assumption we make on the basis is that it is bounded in
the following sense.
\begin{defn}
\label{def:basis} Given a $D$ dimensional subspace $\mathcal{V}\subset\mathcal{C}([0,T],\R)$
and node points $\{c_{j}\}_{j=1}^{D}\subset[0,T]$. For any $\gamma_{\varphi}\geq1$,
we say that a basis $\{\varphi_{j}\}_{j=1}^{D}\subset\mathcal{V}$
is $\gamma_{\varphi}$-bounded if $\varphi_{j}(c_{i})=\delta_{i,j}$
and we have 
\[
\sum_{j=1}^{D}\left|\int_{0}^{t}\varphi_{j}(s)\d s\right|\leq\gamma_{\varphi}T\quad\text{for }t\in[0,T].
\]
\end{defn}

Note that if the constant function $1\in\mathcal{V}$, then we have
\[
1=\sum_{j=1}^{D}1(c_{j})\varphi_{j}(t)=\sum_{j=1}^{D}\varphi_{j}(t).
\]
Hence,
\[
T=\int_{0}^{T}1\d s\leq\sum_{j=1}^{D}\left|\int_{0}^{T}\varphi_{j}(s)\d s\right|\leq\gamma_{\varphi}T.
\]
Therefore, $\gamma_{\varphi}\geq1$ for most elements. Later we will
see that for the space of low degree polynomials or piece-wise low-degree
polynomials, there is a basis on the Chebyshev nodes that is $O(1)$-bounded.

Assuming that $\frac{dx}{dt}$ can be approximated by some element
in $\mathcal{V}$, we have that 
\[
\frac{\d x}{\d t}(t)\sim\sum_{j=1}^{D}\frac{\d x}{\d t}(c_{j})\varphi_{j}(t)=\sum_{j=1}^{D}F(x(c_{j}),c_{j})\varphi_{j}(t).
\]
Integrating both sides, we have 
\begin{equation}
x(t)\approx v+\int_{0}^{t}\sum_{j=1}^{D}F(x(c_{j}),c_{j})\varphi_{j}(s)\d s.\label{eq:approx_fix_point}
\end{equation}
This suggests the following operator from $\mathcal{C}([0,T],\R^{d})$
to $\mathcal{C}([0,T],\R^{d})$: 
\begin{equation}
\mathcal{T}_{\varphi}(x)(t)_{i}=v_{i}+\int_{0}^{t}\sum_{j=1}^{D}F(x(c_{j}),c_{j})_{i}\varphi_{j}(s)\d s\quad\text{for }i\in[d].\label{eq:def_T_d_operator}
\end{equation}
Equation (\ref{eq:approx_fix_point}) can be written as $x\approx\mathcal{T}_{\varphi}(x).$
To find $x$ that satisfies this, naturally, one can apply a fix point
iteration and the resulting algorithm is called the collocation method.\selectlanguage{english}%

